\section{Conclusion}\label{section:conc}

Evidence from equity markets provides resounding support for redistribution-based models of democratization. Democratizations lower stock valuations and raise risk premia substantially across several proxies in data covering 90 countries over 200 years. These results cannot be explained by increased macroeconomic risk nor do other periods of high political or regime transition risk have the same effect. Exogenous variation coming from a change in Catholic church doctrine confirms that risk premia rise with the probability of a successful democratization. 

Redistribution risk can explain these results. In the data, redistribution follows successful democratizations: the size of the public sector and measures of economic competition rise, and income inequality and measures of corruption fall.  Moreover, democratizations with higher redistribution risk see a substantially larger rise in risk premia than other democratizations. A redistribution-based model of democratic transitions with asset prices and incomplete markets can fully explain the results. It can also explain the lack of an asset pricing effect observed in autocratizations.

\label{par:redistribution_highlight}
The analysis highlights several potential channels of redistribution that generate the asset pricing results. That said, in standard macroeconomic models, more redistribution would generally lower growth. This is at odds with empirical evidence that suggests democracy causes higher growth \citep{Acemoglu2019}. Reconciling this disparity would be a natural path for future research. 

The paper also highlights a potential resolution to this apparent contradiction. While much work has focused on declines in inequality or increasing taxation, this paper shows that increased competition and a loss of government consumption for the elites can also play a role in transitions. Indeed, while taxes and inequality are certainly important, they may not play a central role in all democratizations, in particular transitions from left-wing autocracies that are not captured when examining stock market data. My results highlight that approximately half of all redistribution risk impounded in asset prices comes from these other channels. Exploring these channels could help the field to better understand not just periods of democratization, but how policy and political risk affect individual, firm, and government decision making more broadly. It would also illuminate how reduction in barriers to entry relate to economic growth.

Finally, this paper shows that any financial history of the last 200 years that excludes democratizations is incomplete.  In doing this, it provides new avenues of study in consumption-based models by focusing on political institutions and how they interact with the distribution of resources. In a model with incomplete markets, redistribution shocks can have large consequences on asset prices. This means that neither an increase in the probability of a large drop in aggregate consumption nor an increase in the volatility of aggregate consumption is necessary for an increase in risk premia. The consumption risk faced by relatively wealthy investors need only be affected. This paves the road for the risk of redistribution to be a primary historical driver of asset prices.
